\begin{exercise}{42}
Show that if the coherent nerve $\mathrm{N}(\mathcal C)$ of a simplicial category is isomorphic
to the nerve of an ordinary category, then the underlying category $u\mathcal C$ is isomorphic
to the homotopy category $\pi(\mathcal C)$. Make explicit the coherent nerve of the following simplicial
category $\mathbb B^{\mathrm{simp}}(G)$: There is only one object, and the simplicial set of endomorphisms of
this object is given by $\mathrm{N}(G)$, where $G$ is an abelian group (or a commutative monoid, if you wish). Deduce from
the explicit analysis that, for nontrivial $G$, $\mathrm{N}(\mathbb B^{\mathrm{simp}}(G))$ is not isomorphic to the nerve of a category although
$u\mathbb B^{\mathrm{simp}}(G)\cong\pi(\mathbb B^{\mathrm{simp}}(G))$.\footnote{The book says ``a group (of a monoid, if you wish)'';
read ``an abelian group (or a commutative monoid, if you wish)''.
A unital simplicial multiplication on $\mathrm{N}(G)$ forces commutativity by interchange.
Nontriviality is needed only for the non-nerve conclusion: $G=1$ gives $\Delta^0$.
The printed final formula has $u\mathrm{N}(-)$ and $\pi(\mathrm{N}(-))$;
remove these outer nerves, since $u$ and $\pi$ of Definition~1.2.48 take simplicial categories as inputs.}

\begin{proof}
The canonical functor $u\mathcal C\to\pi(\mathcal C)$ is
the identity on objects and sends each vertex of a hom-simplicial
set to its component. An edge $f\to g$ in $\mathrm{Hom}_{\mathcal C}(x,y)$
gives, by Definition~1.2.56, a $2$-simplex of $\mathrm{N}(\mathcal C)$
with vertices $(x,x,y)$, successive edges $\mathrm{id}_x,g$,
and long edge $f$.
In an ordinary nerve its edges satisfy $f=g\,\mathrm{id}_x=g$.
Thus every edge in each hom-simplicial set has equal endpoints,
so each component contains exactly one vertex.
The canonical functor is therefore an isomorphism.

For the example, take an abelian group $G$ and write $BG$
for its one-object category.
The commutativity restriction is necessary: a unital simplicial
multiplication on $\mathrm{N}(BG)$ induces, by full faithfulness
of the ordinary nerve (Corollary~1.2.20), a monoid homomorphism
$m\colon G\times G\to G$ with $m(a,1)=a$ and $m(1,b)=b$.
Since $(a,1)$ and $(1,b)$ commute, $ab=ba$.
Conversely, for abelian $G$, pointwise multiplication on
$\mathrm{N}(BG)_k=G^k$ is simplicial, associative and unital,
and defines $\mathcal B=\mathbb B^{\mathrm{simp}}(G)$.

Here is its coherent nerve in every degree:
\[
 \mathrm{N}(\mathcal B)_n\cong G^{\binom n2}.
\]
The coordinates are $a_{ij}\in G$ for $0\le i<j\le n$
with $j\ge i+2$; put $a_{ii}=a_{i,i+1}=1$.
For a subset $S=\{s_0<\cdots<s_k\}$ containing $i,j$
and contained in $[i,j]$, put
\[
 p_{ij}(S)=\prod_{r=0}^{k-1}a_{s_r,s_{r+1}}.
\]
The hom-component on $\mathrm{N}(P_{i,j})$ is the nerve of
the functor $P_{i,j}\to BG$ sending an inclusion $S\subseteq T$ to
$p_{ij}(S)p_{ij}(T)^{-1}$.
These functors respect enriched composition because the products
$p$ multiply under union of subsets meeting at an endpoint.
Conversely, a hom-functor determines
$p_{ij}(S)$ as the image of $S\subseteq[i,j]$;
composition under union forces this image to be the product
of the $a_{s_r,s_{r+1}}$, where
$a_{ij}=p_{ij}(\{i,j\})$.
Since all arrows of $BG$ are invertible, these potentials
determine every inclusion's image. This proves the bijection.

For every monotone map $\alpha\colon[m]\to[n]$, the simplicial
operator is explicitly
\[
 (\alpha^*a)_{ij}
 =a_{\alpha(i),\alpha(j)}
  \left(\prod_{r=i}^{j-1}a_{\alpha(r),\alpha(r+1)}\right)^{-1}.
\]
Indeed, its coordinate is the image of the inclusion
$\{\alpha(i),\alpha(j)\}\subseteq\alpha([i,j])$.
This gives all faces and degeneracies, including collapsed
intervals by the convention $a_{ii}=1$.

In particular the coherent nerve has one vertex, one edge,
and $G$ as its set of $2$-simplices.
An ordinary category whose nerve has one vertex and one edge
is the terminal category, whose nerve has only one $2$-simplex.
Thus $G\ne1$ gives the required non-nerve.
Nevertheless $u\mathcal B\cong\pi(\mathcal B)\cong[0]$:
both have one object and one morphism, since
$\mathrm{N}(BG)$ has one vertex and one component.
\end{proof}
\end{exercise}
